Digital learning has become an important component of professional training,
particularly in environments where compliance, safety and traceability matter.
An effective learning platform must do more than display documents: it must
connect the training material, assessment, learner progress and validation
decision in one coherent workflow.

The Safyron Learning project addresses this need through a role-oriented web
application. Learners access published modules and their supporting files.
Instructors create the modules, prepare assessments and review submissions.
Administrators oversee users, modules and global activity. The platform also
uses artificial intelligence to reduce the time required to draft questions
from source documents, while preserving human review before publication.

The internship was therefore both a development assignment and an exercise in
software quality. Several business rules had to remain consistent across the
Angular interface, REST API and relational database. A deleted module must not
leave orphaned quizzes or files; an answer must remain associated with the
correct question; and a learner must not obtain final certification before the
instructor has reviewed the submission.

This report follows the structure of the project. Chapter 1 presents the
context, problem statement, stakeholders and scope. Chapter 2 formalizes the
requirements and explains the architecture and design. Chapter 3 presents the
implementation, validation results, live website interface and technical
limitations. The conclusion summarizes the contribution and proposes future
work.
